\documentclass{article}
\usepackage{amsmath,amssymb}
\title{Expectations and Moments}
\author{}

\begin{document}

\maketitle

\section*{Expectations and moments}
The probability distribution defines weighted averages over the sample space, where the weight of each event is equal to its probability. These are called \emph{expected values}.

\subsection*{Discrete case}
For the discrete case,
\[
\mathbb{E}[f(X)] = \sum_{k=1}^n f(x_k)\,p(x_k).
\]

\subsection*{Continuous case}
While for the continuous case,
\[
\mathbb{E}[f] = \int_{-\infty}^{\infty} f(x)\,p(x)\,dx.
\]


\section*{Mean of a distribution}
The mean of the distribution is the expectation of the random variable itself. We denote
\[
\mu \equiv \mathbb{E}[X] = \overline{X} = \langle X \rangle =
\begin{cases}
\displaystyle \sum_k x_k\,p(x_k), & \text{(discrete)}\\[6pt]
\displaystyle \int_{-\infty}^{\infty} x\,p(x)\,dx, & \text{(continuous)}
\end{cases}
\]

In the case of an infinite sample space, whether continuous or discrete, the mean is not guaranteed to exist since the integral or the sum might not converge.

\section*{Moments of a distribution}
The moments of a distribution are the expectation of powers of the random variable itself.
\[
\mu_\ell \equiv \mathbb{E}\left[X^\ell\right] \equiv \left\langle X^\ell \right\rangle = 
\begin{cases} 
\displaystyle\sum_k x_k^\ell p(x_k), & \text{(discrete)}\\[6pt]
\displaystyle\int_{-\infty}^{\infty} x^\ell p(x) \,\mathrm{d}x, & \text{(continuous)}
\end{cases}
\]

If all the moments are known -- and if they exist -- they can be used to get the expectation of other functions using the \textbf{linearity} of the expectation operator:
\begin{align*}
\mathbb{E}[c\,f(X)] &= c\,\mathbb{E}[f(X)], \\
\mathbb{E}[f(X) + g(X)] &= \mathbb{E}[f(X)] + \mathbb{E}[g(X)]
\end{align*}

\section*{Variance and standard deviation}
Of particular interest is the second moment, in combination with the mean, defining the \textbf{variance}
\begin{align*}
\sigma^2 = \operatorname{Var}(X) &\equiv \mathbb{E}\left[(X - \mu)^2\right] \\
&= \mathbb{E}\left[X^2 - 2\mu X + \mu^2\right] \\
&= \mathbb{E}\left[X^2\right] - 2\mu \mathbb{E}[X] + \mu^2 \\
&= \mathbb{E}\left[X^2\right] - \mu^2 \\
&= \mathbb{E}\left[X^2\right] - \mathbb{E}[X]^2
\end{align*}
\section*{Higher moments characterize properties of a distribution}
\textbf{Variance} -- dispersion measure based on the second moment:
\[
\sigma^2 = \mathbb{E}\left[(X-\mu)^2\right] = \int (x-\mu)^2\,p(x)\,dx
\]

\noindent\textbf{Skewness} -- asymmetry parameter based on third moments; dimensionless (normalized cumulant):
\[
s \equiv \frac{\mathbb{E}\left[(X-\mu)^3\right]}{\sigma^3} = \mathbb{E}\left[\left(\frac{X-\mu}{\sigma}\right)^3\right]
\]

\noindent\textbf{Kurtosis} -- measure of tail "weights" in terms of fourth moments; zero for a Gaussian (excess kurtosis), bounded below by -1:
\[
\kappa \equiv \frac{\mathbb{E}\left[(X-\mu)^4\right]}{\sigma^4} - 3
\]

\section*{Covariance and correlation}
For any two random variables, not necessarily independent or identically distributed, their covariance is defined as:
\[
\operatorname{Cov}(X,Y) \equiv \mathbb{E}[(X - \mu_x)(Y - \mu_y)] = \mathbb{E}[XY] - \mu_x\mu_y
\]

\noindent The correlation is proportional to the covariance,
\[
\rho(X,Y) = \operatorname{Corr}(X,Y) \equiv \frac{\operatorname{Cov}(X,Y)}{\sqrt{\operatorname{Var}(X)\operatorname{Var}(Y)}} = \mathbb{E}\left[\left(\frac{X-\mu_x}{\sigma_x}\right)\left(\frac{Y-\mu_y}{\sigma_y}\right)\right]
\]

\noindent Dividing the covariance by the standard deviations makes the correlation a pure number, and
\[
-1 \leq \rho(X,Y) \leq +1
\]

\section*{Common Distributions}

\subsection*{Uniform Distribution}

\begin{itemize}
    \item \textbf{Standard form:}
    \begin{align*}
        p(x) &= \begin{cases}
            1, & x \in [0,1] \\
            0, & \text{otherwise}
        \end{cases}, \\[0.8em]
        \operatorname{Prob}(a < X < b) &= \int_a^b p(x)\,dx = b-a, \\[0.8em]
        F(x) &= \int_{-\infty}^x p(x')\,dx' = x
    \end{align*}

    \item \textbf{Moments:}
    \begin{align*}
        \mu &= \int_{-\infty}^{\infty} x\,p(x)\,dx = \int_0^1 x\,dx = \frac{1}{2}, \\[0.8em]
        \mu_\ell &= \int_0^1 x^\ell\,dx = \frac{1}{\ell + 1}, \\[0.8em]
        \sigma^2 &= \int_0^1 \left(x - \frac{1}{2}\right)^2\,dx = \frac{1}{12}.
    \end{align*}
\end{itemize}

\subsection*{Binomial Distribution}

\begin{itemize}
    \item Identify events by total number of successes versus failures, independent of the order in which they occur.

    \item \textbf{Two parameters:} probability of success $p$, number of trials $n$.
\end{itemize}

\begin{align*}
    f(k;n,p) &= \binom{n}{k}p^k q^{n-k}, \quad \text{where} \\[0.8em]
    q &= 1-p \quad \text{is the probability of failure,} \\[0.8em]
    \binom{n}{k} &= \frac{n!}{k!(n-k)!} = \binom{n}{n-k}.
\end{align*}

\begin{itemize}
    \item \textbf{Mean value (the hard way):}
\end{itemize}

\begin{align*}
    \mu = \mathbb{E}[X]
    &= \sum_{k=0}^{\infty} k f(k;n,p) \\[0.8em]
    &= \sum_{k=1}^{n} \frac{k n!}{k!(n-k)!}p^k q^{n-k} \\[0.8em]
    &= np \sum_{k=1}^{n} \frac{(n-1)!}{(k-1)!(n-k)!}p^{k-1}q^{(n-1)-(k-1)} \\[0.8em]
    &= np \sum_{k'=0}^{n'} f(k';n',p), \quad \text{where } k'=k-1,\ n'=n-1 \\[0.8em]
    &= np.
\end{align*}

\end{document}
